% Adaptación del preámbulo de Paper III v1.5 español; no plantilla standalone de Pandoc.
\documentclass[spanish,11pt,a4paper]{article}
\usepackage[margin=1in]{geometry}
\usepackage{fontspec}
\setmainfont{lmroman10-regular.otf}[BoldFont=lmroman10-bold.otf,ItalicFont=lmroman10-italic.otf,BoldItalicFont=lmroman10-bolditalic.otf]
\setsansfont{lmsans10-regular.otf}[BoldFont=lmsans10-bold.otf]
\setmonofont{lmmono10-regular.otf}[Scale=MatchLowercase]
\newlength{\LeanCharWidth}
\usepackage{amsmath,amssymb}
\usepackage[spanish,es-noshorthands,es-tabla]{babel}
\usepackage{microtype,parskip}
\usepackage{graphicx,float,caption}
\usepackage{longtable,booktabs,array,calc}
\usepackage{xcolor}
\usepackage{xurl}
\usepackage{hyperref,bookmark}
\usepackage{etoolbox,needspace}
\setcounter{secnumdepth}{-\maxdimen}
\setlength{\emergencystretch}{3em}
\providecommand{\tightlist}{\setlength{\itemsep}{0pt}\setlength{\parskip}{0pt}}
\newcounter{none}
\AtBeginEnvironment{longtable}{\Needspace{8\baselineskip}\small}
\pretocmd{\section}{\Needspace{5\baselineskip}}{}{}
\pretocmd{\subsection}{\Needspace{4\baselineskip}}{}{}
\pretocmd{\subsubsection}{\Needspace{4\baselineskip}}{}{}
\captionsetup{font=small,labelfont=bf}
\urlstyle{same}
\hypersetup{colorlinks=true,linkcolor=blue,urlcolor=blue,citecolor=blue,pdflang={es-ES}}
\title{Companion de formalización de la estabilidad cordal exacta}
\author{Juan Pablo Traverso Gianini}
\date{}\begin{document}\raggedbottom\maketitle
Preprint 1.0, 18 de octubre de 2026. Aplicación formal de los Papers IV
y VI. Este companion identifica el paquete final de fuentes, su
enunciado público y sus límites de confianza. La nota desarrolla la
consecuencia de estabilidad enunciada de los papers fuente; no afirma
una nueva solución del Problema 81 de Erdős.

\subsection{Enunciado público y
alcance}\label{enunciado-puxfablico-y-alcance}

La interfaz pública final es
\nolinkurl{ChordalStabilityNote.eventual_root_stability}. Afirma que
existe un umbral de orden N tal que todo grafo G sobre Fin n, con n al
menos N, que satisface cordalidad, \texttt{cp\ G\ \textless{}=\ Mval\ n}
y \texttt{2000\ *\ (Mval\ n\ -\ cp\ G)\ \textless{}=\ n}, satisface
\nolinkurl{RootStability} para el conjunto efectivo definido por grados
\texttt{\{v\ \textbar{}\ n\ \textless{}\ 2\ *\ G.degree\ v\}}.

La cota superior explícita \texttt{cp\ G\ \textless{}=\ Mval\ n} hace
que la resta natural coincida con el déficit no negativo del manuscrito.
Esta interfaz final no supone un testigo estructural, una cota local de
modificación, una fórmula exacta de partición ni una raíz. Todo ello se
obtiene en su demostración. El umbral es el máximo entre el umbral del
redondeo uniforme y \texttt{2\ *\ 10\^{}13}; el término numérico no
reemplaza el umbral de redondeo desconocido.

\nolinkurl{RootStability} es un registro de conclusiones. Incluye:

{\def\LTcaptype{none} % do not increment counter
\begin{longtable}[]{@{}
  >{\raggedright\arraybackslash}p{(\linewidth - 2\tabcolsep) * \real{0.5000}}
  >{\raggedright\arraybackslash}p{(\linewidth - 2\tabcolsep) * \real{0.5000}}@{}}
\toprule\noalign{}
\begin{minipage}[b]{\linewidth}\raggedright
Campo
\end{minipage} & \begin{minipage}[b]{\linewidth}\raggedright
Conclusión matemática
\end{minipage} \\
\midrule\noalign{}
\endhead
\bottomrule\noalign{}
\endlastfoot
\nolinkurl{isClique} & La raíz definida por grados es un clique \\
\nolinkurl{exactFormula} & \texttt{cp\ =\ B(n,p)\ -\ A\ -\ 2m} sobre los
enteros \\
\nolinkurl{four} & \texttt{cp\ =\ cpLe\ 4} \\
\nolinkurl{deficitFormula} & \texttt{M-cp\ =\ (M-B)+A+2m} \\
\nolinkurl{editIdentity} & La distancia de edición a la referencia split
completa es igual a \nolinkurl{A+m} \\
\nolinkurl{editBound} & Esa distancia de edición es a lo sumo
\nolinkurl{M-cp} \\
\nolinkurl{square} & \texttt{(6p-2n-1)\^{}2\ \textless{}=\ 24(M-cp)+25} \\
\nolinkurl{canonical} & La raíz es exactamente el conjunto determinado por
el criterio de grados \\
\end{longtable}
}

La fórmula entera de B evita la resta truncada. La desigualdad
cuadrática se expresa sobre los racionales. Estas elecciones conservan
el significado aritmético del manuscrito.

\subsection{Adaptadores físicos y dependencias de la
demostración}\label{adaptadores-fuxedsicos-y-dependencias-de-la-demostraciuxf3n}

\nolinkurl{PartitionAdapter} convierte la partición exacta efectiva de IV y
la completación de su empaquetamiento físico en la definición de
partición en cliques de VI. Identifica las aristas exteriores, las
aristas del núcleo y los enlaces ausentes efectivos con las funciones de
recuento de VI. Una raíz que es un clique tiene \texttt{Cdef\ =\ 0}.
Estas identificaciones se demuestran; no se introducen como axiomas de
interfaz.

\nolinkurl{WindowArithmetic} y \nolinkurl{GraphWindow} establecen las cotas
escalares de entrada, la comparación entre el grado máximo de
modificación y la distancia total de edición, y el reconocimiento del
núcleo por grados. \nolinkurl{PaidRootAssembly} utiliza entonces una
partición física y sus cuentas pagadas para invocar el teorema
paramétrico efectivo \nolinkurl{PaperVI.cp_window_K4} con parámetro
cuatro. Este teorema intermedio tiene una hipótesis de raíz pagada,
mientras que la interfaz pública final no la tiene.

\nolinkurl{NearAssembly} aplica este resultado intermedio a
\nolinkurl{W.regularized.root} de \nolinkurl{NearStructureWitness} de IV. La
distinción es esencial: las cuentas físicas pertenecen a la raíz
regularizada, no a otro campo denominado core. \nolinkurl{FarGap} excluye
el régimen lejano comparando la partición completada de un
empaquetamiento redondeado efectivo con el óptimo efectivo
\texttt{cp\ G}. \nolinkurl{GlobalAssembly} proporciona el óptimo fraccional
certificado, invoca el redondeo uniforme, obtiene el testigo cercano
mediante localización y ensambla la conclusión.

La cadena de demostración seleccionada requiere
\nolinkurl{uniformTransferAt} de IV,
\nolinkurl{NearH1Localization.exists_nearStructureWitness_of_nearRegime},
su dependencia de localización por descenso, los adaptadores físicos y
la construcción exacta paramétrica de VI. Excluye
\nolinkurl{PaperIV.Erdos81Unconditional}. La aplicación mantiene su
dependencia matemática de IV y VI, aunque tiene su propio proyecto de
fuentes y su propia interfaz final. No utiliza el teorema de
transferencia del Paper V ni los resultados adicionales de recourse de
VI. El cono de entrada de 18 módulos de VI es el subconjunto paramétrico
de fuentes congeladas utilizado en esta demostración; la nota principal
cita el manuscrito editorial vigente sin afirmar que sus módulos
adicionales formen parte de este cono paramétrico de fuentes.

\subsection{Validación final y paquete de
fuentes}\label{validaciuxf3n-final-y-paquete-de-fuentes}

El proyecto utiliza \nolinkurl{leanprover/lean4:v4.35.0}, con Mathlib
fijada en \nolinkurl{56a005855ea17e0b0958f73ec801ae140b0a1679}. Su cierre
de fuentes tiene 343 módulos: ocho módulos de esta nota, 317 del cono de
fuentes de IV y 18 de la construcción paramétrica de VI. Los módulos
heredados de IV conservan su texto fuente y se atribuyen en
\nolinkurl{AUTHORSHIP_NOTICE.md}. Los módulos propios y controles llevan
cabeceras de copyright Apache-2.0; la licencia completa y el aviso de
atribución son \nolinkurl{lean435/LICENSE} y \nolinkurl{lean435/NOTICE}. Los
módulos propios y los controles de la nota identifican a Juan Pablo
Traverso Gianini como autor. Los módulos de VI están sincronizados con
la versión de fuentes con crédito de autor registrada en el manifiesto;
los créditos existentes de terceros permanecen sin cambios.

El build de la raíz canónica comprueba \nolinkurl{ChordalStabilityNote}.
Reutiliza artefactos completos de caché para las dependencias sin
cambios, valida sus trazas y recompila las fuentes modificadas. El
archivo de fuentes es portable y no contiene cachés ni uniones locales
de directorios de desarrollo. Esto no afirma que se hayan recompilado
desde cero todos los módulos previos ni Mathlib.

La auditoría completa comprueba primero que todos los módulos del
manifiesto de fuentes estén presentes en el entorno Lean importado; la
ausencia de un módulo provoca el rechazo.

La auditoría seleccionada de la raíz sigue 7,125 dependencias de
declaraciones, exige los insumos estructurales y de redondeo de IV y el
teorema paramétrico de VI, y excluye
\nolinkurl{PaperIV.Erdos81Unconditional}. La auditoría completa del
proyecto local cubre 7,393 declaraciones y visita 10,976 dependencias.
Ambas permiten únicamente \nolinkurl{propext}, \nolinkurl{Classical.choice} y
\nolinkurl{Quot.sound}; no se acepta ningún axioma matemático adicional ni
\nolinkurl{sorryAx}. \nolinkurl{evidence/BUILD_CHECKS.json} registra los
comandos y códigos de salida finales.
\nolinkurl{FINAL_FORMAL_VALIDATION.json} los vincula con las fuentes y
las comprobaciones de la entrada pública.

\Needspace{8\baselineskip}
El archivo inmutable es
\nolinkurl{NOTE_II_LEAN435_FINAL_SOURCE_v1.0.zip}. Su SHA-256 es:

\begin{verbatim}
b24c5543d719779568ee9c099d453c170d4c69cbd4a31b293f4ee8375e287531
\end{verbatim}

\Needspace{8\baselineskip}
\nolinkurl{SOURCE_MANIFEST.json} registra todos los hashes de fuentes y
configuración. Su SHA-256 es:

\begin{verbatim}
8bcdf425cee099e6e157dffa4356bb720e6fcd47ec138a2d4bac10b710587f75
\end{verbatim}

El archivo \nolinkurl{lean435/OWNERSHIP_AND_LICENSE.json} enumera, con
sus hashes, todas las fuentes Lean propias, incluidos los cortes
reutilizados de la serie, que el autor distribuye bajo Apache 2.0. Las
cabeceras completas nuevas corresponden a los módulos específicos de
esta nota y sus puntos de auditoría; las cabeceras heredadas y las
atribuciones matemáticas se conservan. Las dependencias externas
mantienen sus propias licencias.

\Needspace{12\baselineskip}
\subsection{Reproducción}\label{reproducciuxf3n}

\Needspace{8\baselineskip}
Extraiga el archivo y ejecute:

\begin{verbatim}
cd lean435
lake exe cache get
lake build ChordalStabilityNote
lake env lean -j 1 -s 32768 Endpoint435Gate.lean
lake env lean -j 1 -s 32768 FreezeFullAudit.lean
\end{verbatim}

Conserve el archivo \nolinkurl{lake-manifest.json} suministrado; el comando
de caché obtiene artefactos para las dependencias fijadas sin solicitar
una nueva selección de versiones.

El build de la raíz ejecuta la auditoría seleccionada de dependencias y
axiomas. El último comando audita el conjunto completo de declaraciones
locales importadas. Las dependencias locales están incluidas, de modo
que no se requiere una copia de trabajo de los papers principales.

El resultado final conserva explícitamente
\texttt{cp\ G\ \textless{}=\ Mval\ n}. Sin esta hipótesis, la resta
natural podría convertir un déficit matemático negativo en cero. En el
manuscrito corresponde a la cota inferior del déficit. La cota superior
es \texttt{2000\ *\ (Mval\ n\ -\ cp\ G)\ \textless{}=\ n}; la división
entera no debilita la ventana. El Paper IV proporciona por separado la
cota superior eventual. Combinar su umbral con el de esta nota da el
corolario escrito correspondiente para grafos cordales suficientemente
grandes en la ventana indicada; el resultado congelado conserva la
hipótesis explícita de déficit no negativo.

\subsection{Correspondencia con el manuscrito y
alcance}\label{correspondencia-con-el-manuscrito-y-alcance}

El argumento formal utiliza objetos efectivos en todo momento. Un
empaquetamiento del régimen lejano completado se compara con el óptimo
efectivo. El testigo cercano proporciona una raíz regularizada junto con
su empaquetamiento físico y sus cuentas. Los adaptadores identifican
esas cuentas con las funciones de VI. El teorema de raíz con presupuesto
cubierto demuestra la ventana de entrada y aplica la construcción local
exacta. Cada campo de \nolinkurl{RootStability} es una conclusión, y el
criterio de grados identifica la misma raíz al final.

Las fuentes Markdown \nolinkurl{CHORDAL_STABILITY_NOTE_v1.0_en.md} y
\nolinkurl{CHORDAL_STABILITY_NOTE_v1.0_es.md} conservan el teorema, las
constantes y las ecuaciones expuestas de la deducción establecida. Este
companion describe únicamente el paquete final validado de fuentes. No
afirma un teorema de transferencia inversa, un umbral mejorado, una
nueva extensión de vértices excepcionales ni una asíntota de conteo.

\Needspace{14\baselineskip}\subsection{Asistencia y
responsabilidad}\label{asistencia-y-responsabilidad}

La nota principal conserva la atribución del Paper IV a Claude, de
Anthropic, ChatGPT/Codex, de OpenAI, Aristotle, de Harmonic, Certo y
Jacobian por el trabajo heredado. La contribución de Aristotle a ese
trabajo incluyó la exploración, la búsqueda de contraejemplos, la
elaboración de pruebas y su formalización y revisión; no sólo tradujo
pruebas terminadas. Para esta nota, ChatGPT/Codex ayudó con el
argumento, las interfaces formales, la validación y la presentación;
Aristotle aportó una propuesta separada de formalización. Las
regresiones de aritmética exacta complementan las comprobaciones de
fuentes. Los cálculos se revisaron según su alcance; una búsqueda finita
no reemplaza una prueba universal. La prioridad bibliográfica y la
revisión externa son cuestiones separadas.

\Needspace{3\baselineskip}El autor conserva la responsabilidad por los argumentos, las citas, el
código y la presentación. Ningún sistema de IA figura como autor.

\end{document}
